\documentclass[10pt,twocolumn]{article}
\usepackage[margin=1.8cm]{geometry}
\usepackage{times}
\usepackage{titlesec}
\usepackage{booktabs}
\usepackage{enumitem}
\usepackage{parskip}
\usepackage{amsmath}
\usepackage{graphicx}

\titlespacing{\section}{0pt}{4pt}{2pt}
\titlespacing{\subsection}{0pt}{3pt}{1pt}
\setlist[itemize]{noitemsep, topsep=1pt, leftmargin=*}

\title{\vspace{-1.5em}\textbf{Deep Reinforcement Learning for Automated Stock Trading}\\[2pt]
\normalsize CS5242 Group Project Proposal \textemdash{} Group ID24}
\author{Rutwik Shrikant Dharanappagoudar \quad Zhu Yufan \quad Fei Xin}
\date{\vspace{-1.2em}}

\begin{document}
\maketitle
\thispagestyle{empty}
\vspace{-1em}

%------------------------------------------------------------
\section{Motivation}
Algorithmic trading is a canonical sequential decision-making problem: given
market observations, an agent must decide when to \emph{buy}, \emph{sell},
or \emph{hold} a position to maximise long-term profit.
Classical rule-based strategies such as moving-average crossovers are rigid
and fail to adapt to the non-stationary nature of financial markets.
Deep Reinforcement Learning (DRL) offers a principled, data-driven framework
for learning trading policies end-to-end, without hand-crafted rules,
making it a natural and compelling application of the deep learning
concepts studied in this course.

%------------------------------------------------------------
\section{Project Description}
We design and implement a complete stock trading system \textbf{from scratch}
in PyTorch, where an RL agent learns to trade a selection of
Dow Jones Industrial Average (DJIA) stocks purely from historical price data.

The agent operates inside a custom \texttt{TradingEnv}, a Gym-style
environment we build ourselves that simulates a realistic single-asset
trading scenario: at each daily time step the agent observes the current
market state, selects an action, and receives a reward signal based on
the resulting profit or loss.
The environment explicitly models transaction costs to discourage excessive
trading and encourage stable, profitable policies.

We implement and compare two DRL algorithms of increasing sophistication:
\begin{itemize}
  \item \textbf{Deep Q-Network (DQN):} combines Q-learning with a neural
        network function approximator, an experience replay buffer for
        sample decorrelation, and a periodically updated target network
        for training stability.
  \item \textbf{Double DQN (DDQN):} extends DQN by decoupling action
        selection from value estimation, reducing the systematic
        overestimation of Q-values that degrades DQN performance.
\end{itemize}
Both agents share the same environment and are benchmarked against a
passive \emph{buy-and-hold} strategy to assess whether the learned
policies genuinely add value.

%------------------------------------------------------------
\section{Proposed Solution}

\textbf{Dataset.}
Daily OHLCV (Open, High, Low, Close, Volume) data for selected DJIA
stocks from 2009 to 2024, fetched via \texttt{yfinance}.
The data is split chronologically: 70\% training, 15\% validation,
15\% testing, to avoid look-ahead bias.

\textbf{State space.}
At each time step the agent receives a feature vector encoding:
(1) normalised daily return $\frac{p_t - p_{t-1}}{p_{t-1}}$,
(2) price deviation from the 5-day moving average,
(3) price deviation from the 20-day moving average,
(4) 14-day Relative Strength Index (RSI-14), and
(5) the agent's current position $\in\{-1,\,0,\,1\}$.
All features are normalised to zero mean and unit variance.

\textbf{Action space.}
Three discrete actions: Buy\,(+1), Hold\,(0), Sell\,($-1$).

\textbf{Reward function.}
\[
  r_t \;=\; \mathrm{pos}_t \times \frac{p_t - p_{t-1}}{p_{t-1}}
           \;-\; c\cdot\mathbf{1}[\text{trade executed}]
\]
where $c = 0.001$ penalises each executed trade to model transaction costs.

\textbf{Network architecture.}
A 3-layer MLP: input $\to$ 128 $\to$ 64 $\to$ \#actions, with ReLU
activations and a linear output layer. Trained with the Adam optimiser
($\mathrm{lr} = 10^{-3}$, batch size 64), replay buffer of 10{,}000
transitions, and target network synchronised every 100 gradient steps.

\textbf{Evaluation.}
Performance on the held-out test set is measured by cumulative portfolio
return, annualised Sharpe ratio, and maximum drawdown, all compared
against the buy-and-hold baseline.

%------------------------------------------------------------
\section{Project Milestones}

\begin{center}
\small
\begin{tabular}{@{}cl@{}}
\toprule
\textbf{Week} & \textbf{Milestone} \\
\midrule
8  & Data pipeline: download, clean, feature engineering \\
9  & \texttt{TradingEnv} + DQN implementation from scratch \\
10 & DDQN; single-stock (AAPL) training \& debugging \\
11 & Multi-stock experiments; TA Zoom meeting \\
12 & Evaluation metrics, visualisation, baseline comparison \\
13 & Report writing \& slide preparation \\
14 & Video recording; final submission (Apr 26) \\
\bottomrule
\end{tabular}
\end{center}

%=============================================================
% PAGE 2: TEAM CONTRACT (not counted in 1-page limit)
%=============================================================
\newpage
\onecolumn
\begin{center}
  {\LARGE\textbf{Team Contract --- Group ID24}}\\[6pt]
  \textit{CS5242: Neural Networks and Deep Learning, Semester 2 2025/26}
\end{center}

\bigskip
\noindent\textbf{Project Title:} Deep Reinforcement Learning for Automated Stock Trading

\bigskip
\noindent\textbf{Agreed Principles:}
\begin{itemize}
  \item All members commit equal effort and attitude toward the project's success.
  \item Weekly team sync every Monday via Zoom/Teams.
  \item Weekly progress update sent to the assigned TA every Friday by 6\,pm (Weeks 8--13).
  \item Decisions are made by consensus; conflicts resolved via majority vote.
  \item All code is version-controlled in a shared GitHub repository.
\end{itemize}

\bigskip
\noindent\textbf{Task Distribution:}

\bigskip
\begin{center}
\begin{tabular}{p{5.5cm} p{9.5cm}}
\toprule
\textbf{Team Member} & \textbf{Responsibilities} \\
\midrule
Rutwik Shrikant Dharanappagoudar &
  DQN \& DDQN implementation from scratch in PyTorch;
  experience replay buffer; target network update schedule;
  hyperparameter tuning. \\[6pt]
Zhu Yufan &
  Data collection \& preprocessing pipeline (\texttt{yfinance});
  feature engineering (RSI, moving averages, normalisation);
  \texttt{TradingEnv} (Gym-style) environment design and testing. \\[6pt]
Fei Xin &
  Evaluation framework (Sharpe ratio, max drawdown, cumulative return);
  visualisation of training curves and portfolio performance;
  project report and presentation slides. \\[4pt]
\midrule
\textbf{All members} &
  Final video presentation (each member presents their own contribution). \\
\bottomrule
\end{tabular}
\end{center}

\vspace{2em}
\noindent\textbf{Signatures:}

\vspace{2.5em}
\noindent
\begin{tabular}{p{7cm} p{1cm} p{7cm}}
\underline{\hspace{7cm}} & & \underline{\includegraphics[height=1.2cm]{YF_signature.png}\hspace{2.8cm}} \\[3pt]
Rutwik Shrikant Dharanappagoudar \hfill Date: \underline{\hspace{2.5cm}} & &
Zhu Yufan \hfill Date: 08/03/2026 \\[3em]
\underline{\hspace{7cm}} & & \\[3pt]
Fei Xin \hfill Date: \underline{\hspace{2.5cm}} & & \\
\end{tabular}

\end{document}
